% GENERATED FILE. Edit canonical lessons, then run npm run generate:books.
% canonical-source-hash: fnv1a64:f8d8eb3fcca2730c
% canonical-lessons: SA-C362-anusandhanam, SA-C362-prayogah, SA-C362-anvesanam, SA-C362-aveksate, SA-C362-udbhavayati

\chapter{Research}
\label{ch:sa-research}
\hlchaptermodality{362}

\begin{chapteropening}
\textbf{By the end of this chapter:} \emph{I can say what a researcher is working on.}

Complete the last lesson of chapter 362: I can say what a researcher is working on.
\end{chapteropening}

\section[\texorpdfstring{\sk{अनुसन्धानम्} (anusandhānam)}{अनुसन्धानम् (anusandhānam)}]{\sk{अनुसन्धानम्} (anusandhānam) — research; investigation}

\label{lesson:SA-C362-anusandhanam}



\begin{quote}
Before the new one: say the Sanskrit for ungrateful, then the Sanskrit for gives one's word.
\end{quote}

\subsection*{You'll want to know: \sk{अनुसन्धानम्}}
\textbf{\sk{अनुसन्धानम्}} — \emph{anusandhānam} — \textquotedblleft{}research; investigation\textquotedblright{}.

\begin{quote}
\sk{वैज्ञानिकः} \sk{रोगस्य} \sk{अनुसन्धानं} \sk{करोति।}
\end{quote}

\begin{quote}
\emph{vaijñānikaḥ rogasya anusandhānaṃ karoti.}
\end{quote}

\begin{quote}
The scientist is doing research on the disease.
\end{quote}

A neuter noun.

\begin{cousinweb}
\textbf{\sk{अनुसन्धानम्}} (\emph{anusandhānam}) is built on \emph{anu}, \emph{sam} and the root √dhā, \textquotedblleft{}to put\textquotedblright{}; MW gives \emph{anu-saṃdhāna} \textquotedblleft{}investigation, inquiry, searching into, close inspection\textquotedblright{} (MW).
\end{cousinweb}

\begin{grammarlens}[title={What you've built}]
One more word for giving reasons and opinions.
\end{grammarlens}

\subsection*{Your turn}
\begin{itemize}
\raggedright
  \item \emph{Say it:} \emph{anusandhānam}
  \item \emph{Say it:} \emph{anusandhānam}, once more
  \item \emph{Say it:} \emph{vaijñānikaḥ rogasya anusandhānaṃ karoti.} — The scientist is doing research on the disease.
  \item \emph{Recall:} say the Sanskrit for ungrateful, then the Sanskrit for gives one's word, then say \emph{anusandhānam} again
\end{itemize}

\begin{tcolorbox}[breakable,colback=teal!4,colframe=teal!35!black,title={Before you move on}]
What does \sk{अनुसन्धानम्} mean? (\textquotedblleft{}Research; investigation\textquotedblright{}.) Say it once more.
\end{tcolorbox}
\section[\texorpdfstring{\sk{प्रयोगः} (prayogaḥ)}{प्रयोगः (prayogaḥ)}]{\sk{प्रयोगः} (prayogaḥ) — an experiment; a use}

\label{lesson:SA-C362-prayogah}



\begin{quote}
Before the new one: say the Sanskrit for gives one's word, then the Sanskrit for research.
\end{quote}

\subsection*{You'll want to know: \sk{प्रयोगः}}
\textbf{\sk{प्रयोगः}} — \emph{prayogaḥ} — \textquotedblleft{}an experiment; a use\textquotedblright{}.

\begin{quote}
\sk{वैज्ञानिकः} \sk{नूतनं} \sk{प्रयोगं} \sk{करोति।}
\end{quote}

\begin{quote}
\emph{vaijñānikaḥ nūtanaṃ prayogaṃ karoti.}
\end{quote}

\begin{quote}
The scientist carries out a new experiment.
\end{quote}

A masculine noun.

\begin{cousinweb}
\textbf{\sk{प्रयोगः}} (\emph{prayogaḥ}) is built on \emph{pra} and the root √yuj, \textquotedblleft{}to join\textquotedblright{}; MW gives \emph{pra-yoga} \textquotedblleft{}application, use\textquotedblright{} and \textquotedblleft{}practice, experiment\textquotedblright{} (MW).
\end{cousinweb}

\begin{grammarlens}[title={What you've built}]
One more word for giving reasons and opinions.
\end{grammarlens}

\subsection*{Your turn}
\begin{itemize}
\raggedright
  \item \emph{Say it:} \emph{prayogaḥ}
  \item \emph{Say it:} \emph{prayogaḥ}, once more
  \item \emph{Say it:} \emph{vaijñānikaḥ nūtanaṃ prayogaṃ karoti.} — The scientist carries out a new experiment.
  \item \emph{Recall:} say the Sanskrit for gives one's word, then the Sanskrit for research, then say \emph{prayogaḥ} again
\end{itemize}

\begin{tcolorbox}[breakable,colback=teal!4,colframe=teal!35!black,title={Before you move on}]
What does \sk{प्रयोगः} mean? (\textquotedblleft{}An experiment; a use\textquotedblright{}.) Say it once more.
\end{tcolorbox}
\section[\texorpdfstring{\sk{अन्वेषणम्} (anveṣaṇam)}{अन्वेषणम् (anveṣaṇam)}]{\sk{अन्वेषणम्} (anveṣaṇam) — an inquiry; a search}

\label{lesson:SA-C362-anvesanam}



\begin{quote}
Before the new one: say the Sanskrit for research, then the Sanskrit for an experiment.
\end{quote}

\subsection*{You'll want to know: \sk{अन्वेषणम्}}
\textbf{\sk{अन्वेषणम्}} — \emph{anveṣaṇam} — \textquotedblleft{}an inquiry; a search\textquotedblright{}.

\begin{quote}
\sk{आरक्षकः} \sk{चौरस्य} \sk{अन्वेषणं} \sk{करोति।}
\end{quote}

\begin{quote}
\emph{ārakṣakaḥ caurasya anveṣaṇaṃ karoti.}
\end{quote}

\begin{quote}
The policeman searches for the thief.
\end{quote}

A neuter noun.

\begin{cousinweb}
\textbf{\sk{अन्वेषणम्}} (\emph{anveṣaṇam}) is built on \emph{anu}, \textquotedblleft{}after\textquotedblright{}, and the root √iṣ, \textquotedblleft{}to seek\textquotedblright{}; MW gives \emph{anv-√iṣ} \textquotedblleft{}to seek, to search, to aim at\textquotedblright{}, and \emph{anv-eṣaṇa} the seeking itself (MW).
\end{cousinweb}

\begin{grammarlens}[title={What you've built}]
One more word for giving reasons and opinions.
\end{grammarlens}

\subsection*{Your turn}
\begin{itemize}
\raggedright
  \item \emph{Say it:} \emph{anveṣaṇam}
  \item \emph{Say it:} \emph{anveṣaṇam}, once more
  \item \emph{Say it:} \emph{ārakṣakaḥ caurasya anveṣaṇaṃ karoti.} — The policeman searches for the thief.
  \item \emph{Recall:} say the Sanskrit for research, then the Sanskrit for an experiment, then say \emph{anveṣaṇam} again
\end{itemize}

\begin{tcolorbox}[breakable,colback=teal!4,colframe=teal!35!black,title={Before you move on}]
What does \sk{अन्वेषणम्} mean? (\textquotedblleft{}An inquiry; a search\textquotedblright{}.) Say it once more.
\end{tcolorbox}
\section[\texorpdfstring{\sk{अवेक्षते} (avekṣate)}{अवेक्षते (avekṣate)}]{\sk{अवेक्षते} (avekṣate) — inspects; looks into}

\label{lesson:SA-C362-aveksate}



\begin{quote}
Before the new one: say the Sanskrit for an experiment, then the Sanskrit for an inquiry.
\end{quote}

\subsection*{You'll want to know: \sk{अवेक्षते}}
\textbf{\sk{अवेक्षते}} — \emph{avekṣate} — \textquotedblleft{}inspects; looks into\textquotedblright{}.

\begin{quote}
\sk{अधिकारी} \sk{कार्यालयम्} \sk{अवेक्षते।}
\end{quote}

\begin{quote}
\emph{adhikārī kāryālayam avekṣate.}
\end{quote}

\begin{quote}
The official inspects the office.
\end{quote}

Past (he …): \sk{सः} \sk{कार्यम्} \sk{अवेक्षितवान्}, \emph{saḥ kāryam avekṣitavān}.

\begin{cousinweb}
\textbf{\sk{अवेक्षते}} (\emph{avekṣate}) is built on \emph{ava}, \textquotedblleft{}down\textquotedblright{}, and the root √īkṣ, \textquotedblleft{}to see\textquotedblright{}; MW gives \emph{avekṣ} (\emph{ava} with √īkṣ) \textquotedblleft{}to perceive, observe\textquotedblright{} and \textquotedblleft{}to take into consideration\textquotedblright{}; KRM gives \emph{īkṣita}, \emph{īkṣitavat} for the root, so \emph{avekṣita} (MW).
\end{cousinweb}

\begin{grammarlens}[title={What you've built}]
One more word for giving reasons and opinions.
\end{grammarlens}

\subsection*{Your turn}
\begin{itemize}
\raggedright
  \item \emph{Say it:} \emph{avekṣate}
  \item \emph{Say it:} \emph{avekṣate}, once more
  \item \emph{Say it:} \emph{adhikārī kāryālayam avekṣate.} — The official inspects the office.
  \item \emph{Recall:} say the Sanskrit for an experiment, then the Sanskrit for an inquiry, then say \emph{avekṣate} again
\end{itemize}

\begin{tcolorbox}[breakable,colback=teal!4,colframe=teal!35!black,title={Before you move on}]
What does \sk{अवेक्षते} mean? (\textquotedblleft{}Inspects; looks into\textquotedblright{}.) Say it once more.
\end{tcolorbox}
\section[\texorpdfstring{\sk{उद्भावयति} (udbhāvayati)}{उद्भावयति (udbhāvayati)}]{\sk{उद्भावयति} (udbhāvayati) — devises; brings into being}

\label{lesson:SA-C362-udbhavayati}



\begin{quote}
Before the new one: say the Sanskrit for an inquiry, then the Sanskrit for inspects.
\end{quote}

\subsection*{You'll want to know: \sk{उद्भावयति}}
\textbf{\sk{उद्भावयति}} — \emph{udbhāvayati} — \textquotedblleft{}devises; brings into being\textquotedblright{}.

\begin{quote}
\sk{वैज्ञानिकः} \sk{नूतनं} \sk{यन्त्रम्} \sk{उद्भावयति।}
\end{quote}

\begin{quote}
\emph{vaijñānikaḥ nūtanaṃ yantram udbhāvayati.}
\end{quote}

\begin{quote}
The scientist devises a new machine.
\end{quote}

Past (he …): \sk{सः} \sk{नूतनम्} \sk{उपायम्} \sk{उद्भावितवान्}, \emph{saḥ nūtanam upāyam udbhāvitavān}.

\begin{cousinweb}
\textbf{\sk{उद्भावयति}} (\emph{udbhāvayati}) is built on \emph{ut}, \textquotedblleft{}up\textquotedblright{}, and the root √bhū, \textquotedblleft{}to be\textquotedblright{}, in its causative form, \textquotedblleft{}to cause to exist\textquotedblright{}; MW gives \emph{ud-bhāvita} \textquotedblleft{}caused to exist, created, produced\textquotedblright{} (MW).
\end{cousinweb}

\begin{grammarlens}[title={What you've built}]
One more word for giving reasons and opinions.
\end{grammarlens}

\subsection*{Your turn}
\begin{itemize}
\raggedright
  \item \emph{Say it:} \emph{udbhāvayati}
  \item \emph{Say it:} \emph{udbhāvayati}, once more
  \item \emph{Say it:} \emph{vaijñānikaḥ nūtanaṃ yantram udbhāvayati.} — The scientist devises a new machine.
  \item \emph{Recall:} say the Sanskrit for an inquiry, then the Sanskrit for inspects, then say \emph{udbhāvayati} again
\end{itemize}

\begin{tcolorbox}[breakable,colback=teal!4,colframe=teal!35!black,title={Before you move on}]
What does \sk{उद्भावयति} mean? (\textquotedblleft{}Devises; brings into being\textquotedblright{}.) Say it once more.
\end{tcolorbox}
